\subsection{Moving the sole falsity guard to a register that is about to be overwritten}
\label{sec:guard-relocation}
The logical proof search need not test falsity at the same command boundary
as earlier candidates.  This gives several command interfaces while retaining
the accepted numerical calculus and its exact full AX bundle.  The point is
to choose a guard whose register handling can share work with the surrounding
source macro.

Use the work tuple $(t,w,d,h)$ and Cantor stack notation $[x;w]=P(x,w)$.
The macro names in this subsection are local to the chosen alternative
interface. Let $A_0(c;x,y)$ denote the predecessor AX command: with
$U=d$, $A=d+x$, $B=h+y$, $C=c$, it returns the full gate and
\[
 [R_1,\ldots,R_7,L(A,B,C,U),Q_Uh,t;w],
 \qquad d=h=0.
\]
Its gate and first eight new heads have tag zero; the ninth has tag two.
The enabled gate is itself an implication code, since it is the literal
negation $N(V_0\to J_0)$.  In particular none of these nine heads or the
gate equals native falsity~$1$.

\paragraph{A top guard on transfer gives a total proof compiler.}
Define
\[
 S_T(t,w,d,h)=
 \begin{cases}
  \mathsf{halt},&t=1,\\
  (h,\operatorname{tail}(w),d,\operatorname{head}(w)),&t\ne1.
 \end{cases}
\]
Keep AX and MP unchanged.  Put
\[
 D=\mathsf{AX}(0;0,0),\qquad
 H=D;S_T^{11},\qquad
 \mathsf{PUSH}=D;S_T^9;\mathsf{MP}(0).
\]
For every natural input these macros terminate normally and satisfy
\[
 H(t,w,d,h)=(t,\operatorname{tail}(w),0,\operatorname{head}(w)),
 \qquad
 \mathsf{PUSH}(t,w,d,h)=(0,[t;w],0,0).
\]
Indeed, the first eleven transfer guards after AX inspect only its gate,
zero, its first eight heads, and its ninth head.  The saved incoming top
appears as output of the eleventh shift and is never tested.  Thus $H$
returns even when $t=1$.  After nine shifts PUSH has a tag-two cache;
its final MP necessarily fails against a tag-zero implication.

The reset macro
\[
 \mathsf{DISCARD}=\mathsf{MP}(1);\mathsf{AX}(1;0,0);
 S_T^{10};\mathsf{MP}(0)
\]
returns $(0,w,0,0)$ from every $(t,w,0,0)$.  Define $J(0)$ to be empty
and $J(u)=\mathsf{MP}(u)$ for $u>0$ on a cleared input.  These give a
total, exact compiler $B_X(t,w,0,0)=(X,w,0,0)$ for every native theorem $X$.
The blocks are as follows, with $M=I(A,B)$:
\begin{align*}
 B_0&=\mathsf{DISCARD},\\
 B_{R_j(a,b,c)}&=\mathsf{DISCARD};\mathsf{AX}(c;a,b);
 S_T^{j+1};H^{9-j}, &&1\le j\le7,\\
 B_{L(a,b,c,u)}&=\mathsf{DISCARD};J(u);\mathsf{AX}(c;a-u,b);
 S_T^9;H, &&u\le a,\\
 B_{\operatorname{Bundle}(a,b,c)}&=
 \mathsf{DISCARD};\mathsf{AX}(c;a,b);H^{10},\\
 B_{Q_iX}&=B_X;\mathsf{PUSH};J(i);S_T;
 \mathsf{AX}(i;0,0);S_T^{10},\\
 B_B&=B_M;\mathsf{PUSH};B_A;\mathsf{PUSH};S_T;S_T;
 \mathsf{MP}(B);\mathsf{AX}(B;0,0);S_T^{11}.
\end{align*}
Each axiom block selects the indicated head and removes exactly the
remaining generated heads through the saved zero.  In GEN the lone shift
tests top zero and loads $X$ into cache immediately before AX.  In MP the
two shifts both test top zero and leave $(A,w,0,M)$.  The literal MP succeeds,
pushing one zero.  The cleanup AX has zero gate and its eleven shifts
remove precisely that zero.  Its last guard tests $Q_B0$, so even $B=1$
returns normally.  The alias $I(0,0)=0$ is handled by this same successful
MP calculation.  Structural induction on the finite proof establishes the
total compiler without a theoremhood condition on incoming $t,w$.
The final word $B_1;S_T$ therefore halts.

\paragraph{A guard at AX entry gives the usual strong compiler.}
Alternatively make transfer total,
\[
 S_0(t,w,d,h)=(h,\operatorname{tail}(w),d,\operatorname{head}(w)),
\]
and define $\mathsf{AX}^{\dagger}$ to halt exactly when incoming $t=1$,
otherwise performing exactly $A_0$.  Keep MP unchanged.  The displayed
compiler blocks, with $S_0$ and $\mathsf{AX}^{\dagger}$ substituted, now
obey the strong contract: they halt or return exactly $(X,w,0,0)$.

Here $H$ and PUSH halt precisely on incoming top one.  DISCARD remains
total because its first MP clears the top before any guard.  All axiom
blocks enter AX with top zero or a selected implication code.  GEN may
halt in PUSH if $X=1$, which satisfies the strong contract.  MP may halt
in its second PUSH if $A=1$, or at cleanup AX if $B=1$; otherwise every
displayed state equality holds.  These are all new halting cases, and the
old numerical zero aliases require no additional exception.  The fixed
suffix $B_1;\mathsf{AX}^{\dagger}(0;0,0)$ halts from every top/stack pair
with counter/cache zero.

\paragraph{The same guard may instead be placed at MP entry.}
With total $S_0$ and unguarded AX, the displayed blocks terminate normally
with their exact outputs.  Add a test of incoming $t=1$ to the entry of
just one chosen command.  Until that test halts, execution follows exactly
the total command computation; thus every block now either halts or
returns its former exact output.  Actual-run theoremhood makes every new
halt sound.  Appending that guarded command to $B_1$ gives completeness.
In particular one may guard MP instead of AX and use the final suffix
$B_1;\mathsf{MP}^{\dagger}(0)$.  This is useful when the tested top is
already being copied or cleared inside the MP implementation.  It does
not require inserting guards into both commands.

\paragraph{Soundness, protected tails, and cofinal completeness.}
On safe runs the tested top is always a native theorem.  Both interfaces
therefore halt only with a native derivation of~1.  Their normal command
outputs preserve theoremhood.  The protected-tail balance proof also
applies unchanged: the three command costs remain $-1,0,+10$, and each
returning transfer removes exactly one explicit theorem head.  In the
second interface an AX halt is already sound, while each returning AX
still adds ten heads.  Thus the existing all-natural credit theorem,
including $d(1)=2$, applies to both interfaces.

The fields remain $[0]$, $[1,x]$, and $[c+2,y,x]$.
If $R$ is the encoded residual for the corresponding fixed refutation
suffix, each candidate
\[
 N_Y=P\bigl(2,P(Y,P(1,R))\bigr)
\]
first executes $\mathsf{AX}(0;1,Y)$ (or its guarded version), then the
suffix.  It either halts on that first call or resets counter/cache and
executes a successful compiled refutation.  Since $N_Y\ge Y$, successful
candidates occur beyond every finite startup cutoff.  The unchanged
scheduler's positive residual strictly decreases after decrement and field
decoding, so it eventually reaches such a candidate unless it has halted
earlier.  Under the inherited native-calculus calibration, either source
interface has halting equivalent to inconsistency of ZFC.

These are complete command-level proofs.  The source implementation must
still prove its exact guard and normal output contracts, scratch
restoration, and wire decoding.  Physical allocation, all-tape simulation,
and blank initialization are separate obligations; a smaller source DAG
or finite sample alone does not discharge them.

\paragraph{An independently chosen generalization binder.}
A further compatible change is to replace only AX's ninth head $Q_Uh$ by
$Q_Ch$, keeping all eight logical rows and the saved top unchanged.
Soundness still follows from unrestricted theorem generalization.
H and PUSH require only the ninth head's tag two, so their contracts and
all axiom and MP blocks remain valid.  Writing $S$ for the transfer operation of the chosen alternative,
the GEN block simplifies to
\[
 B_{Q_iX}=B_X;\mathsf{PUSH};S;\mathsf{AX}(i;0,0);S^{10}.
\]
Before AX the state is $(0,w,0,X)$, so the ninth head is exactly $Q_iX$
and the tenth is zero.  The old counter-setting block $J(i)$ is omitted.
Here the AX gate need not vanish, because $A=0$ and $B=X$; it remains an
implication code and the first transfer safely overwrites it.  This detail
also makes the construction valid with the top-guard transfer.
